\section{Заключение}

В ходе выполнения лабораторной работы у меня возникали некоторые вопросы, на которые я смог найти ответ самостоятельно или в различных источниках. Поэтому приведу ответы на них здесь в качестве результата.

\paragraph{Почему мы не стали «честно» решать СЛАУ?}

Рассмотреть численный метод решения крамера, сказать, что обусловленность высокая

Прийти к прямому методу решения СЛАУ

крч всё по Петрову и Лобанову

Достоинства:
• Может быть применен в любом случае, если у системы есть решение
и на главной диагонали нет 0 элементов.
• Нет дополнительных подготовительных вычислений.
• Результаты могут быть получены за конечное количество операций.
• Низкая погрешность метода, если производить расчёты в последний
момент (например, при вычислении вручную).
Недостатки:
• При вычислении на компьютере образуется значительная набегаю-
щая погрешность за счёт ограниченности разрядной сетки, особенно
при работе с разреженными матрицами (коэффициенты близки к 0).
• Ошибки накапливаются, т. к. расчёт производится последовательно.
• Может занимать существенный объем памяти для больших матриц,
так как всю матрицу нужно держать в памяти.
• Алгоритмическая сложность метода O n3 , где n – число неизвест-
ных.

Сравнить с другим методом из той же категории!!! 